% This document was generated by an LLM.

\documentclass{essay}

\title{The Picard--Lindel\"of Theorem}
\subtitle{Existence and Uniqueness for Initial Value Problems}
\author{}
\date{}

\begin{document}

\maketitle

\section{Introduction: A Question Easy to Overlook}

After learning a handful of techniques for solving ordinary differential equations---separation of variables, integrating factors, exact equations, series methods---it is natural to develop the habit of assuming that every reasonable-looking initial value problem \emph{has} a solution, and that the solution we find by whatever technique is \emph{the} solution. The Picard--Lindel\"of theorem is the result that tells us when these assumptions are actually justified. It answers two distinct questions about the initial value problem
\begin{equation}
  y' = f(x, y), \qquad y(x_0) = y_0:
  \label{eq:ivp}
\end{equation}
\begin{enumerate}[label=(\roman*)]
  \item \textbf{Existence:} does a solution exist at all, at least on some interval around $x_0$?
  \item \textbf{Uniqueness:} if a solution exists, is it the only one?
\end{enumerate}
Both questions have negative answers in general, and the counterexamples are not exotic. The theorem identifies a simple, checkable condition on $f$ under which both answers become ``yes.'' Along the way, its proof introduces a constructive method---Picard iteration---that connects differential equations, integral equations, and power series in a satisfying way.

\section{Prerequisite: The Lipschitz Condition}

The key hypothesis of the theorem is a regularity condition on $f$ that is stronger than continuity but weaker than differentiability. It deserves careful attention, because it is precisely the ingredient that makes uniqueness work.

\begin{definition}[Lipschitz condition in one variable]
  A function $g \colon I \to \R$ defined on an interval $I$ satisfies a \emph{Lipschitz condition} (or ``is Lipschitz'') on $I$ if there exists a constant $L \geq 0$ such that
  \[
    |g(u) - g(v)| \leq L\,|u - v| \qquad \text{for all } u, v \in I.
  \]
  The smallest such $L$ is called the \emph{Lipschitz constant} of $g$ on $I$.
\end{definition}

Geometrically, the condition says that the slope of every secant line through two points of the graph of $g$ is bounded in absolute value by $L$:
\[
  \left| \frac{g(u) - g(v)}{u - v} \right| \leq L.
\]
The graph cannot become arbitrarily steep anywhere on $I$. This is a \emph{global} bound on steepness over the whole interval, which is what distinguishes it from mere differentiability at each point.

\subsection{How Lipschitz relates to continuity and differentiability}

The Lipschitz condition sits strictly between continuity and having a bounded derivative:
\begin{itemize}
  \item \textbf{Lipschitz $\Rightarrow$ continuous.} If $|g(u)-g(v)| \leq L|u-v|$, then making $u$ close to $v$ forces $g(u)$ close to $g(v)$. (In fact Lipschitz functions are \emph{uniformly} continuous.)
  \item \textbf{Bounded derivative $\Rightarrow$ Lipschitz.} If $g$ is differentiable on $I$ and $|g'| \leq L$ throughout $I$, the Mean Value Theorem gives, for any $u, v \in I$,
    \[
      g(u) - g(v) = g'(\xi)(u - v) \quad \text{for some } \xi \text{ between } u \text{ and } v,
    \]
    hence $|g(u)-g(v)| \leq L|u-v|$. This is the practical test used most often: \emph{to check the Lipschitz condition, bound the derivative}.
  \item Neither implication reverses, as the examples below show.
\end{itemize}

\begin{example}[Lipschitz but not differentiable]
  $g(u) = |u|$ on $\R$. The triangle inequality gives $\bigl||u|-|v|\bigr| \leq |u-v|$, so $g$ is Lipschitz with $L = 1$, yet $g$ is not differentiable at $0$. Lipschitz functions may have corners; what they may not have is unbounded steepness.
\end{example}

\begin{example}[Continuous but not Lipschitz]
  $g(u) = \sqrt{|u|}$ on $[-1, 1]$. Take $v = 0$: then
  \[
    \frac{|g(u) - g(0)|}{|u - 0|} = \frac{\sqrt{|u|}}{|u|} = \frac{1}{\sqrt{|u|}} \to \infty \quad \text{as } u \to 0.
  \]
  No finite $L$ can bound all these secant slopes, so $g$ is not Lipschitz on any interval containing $0$. The graph has a vertical tangent there. The same argument applies to $g(u) = u^{1/3}$, which will reappear as the villain of Section~\ref{sec:nonuniqueness}.
\end{example}

\begin{example}[Lipschitz locally but not globally]
  $g(u) = u^2$ on $\R$. On any bounded interval $[-M, M]$ we have $|g'(u)| = |2u| \leq 2M$, so $g$ is Lipschitz there with $L = 2M$. But on all of $\R$ the derivative is unbounded, and no single constant works. Many important right-hand sides in ODE theory (such as $f(y) = y^2$) are of this type: \emph{locally} Lipschitz but not globally. As we will see, this is exactly why the theorem's conclusion is local.
\end{example}

\subsection{Lipschitz in the second variable}

In the initial value problem~\eqref{eq:ivp}, the right-hand side $f(x, y)$ is a function of two variables. The condition we need concerns only its behaviour in the $y$-direction:

\begin{definition}[Lipschitz condition in $y$]
  Let $R$ be a rectangle $\{(x,y) : |x - x_0| \leq a,\ |y - y_0| \leq b\}$ around the initial point. We say $f(x,y)$ is \emph{Lipschitz in $y$ on $R$} if there exists $L \geq 0$ such that
  \[
    |f(x, y_1) - f(x, y_2)| \leq L\,|y_1 - y_2|
    \qquad \text{for all } (x, y_1), (x, y_2) \in R.
  \]
\end{definition}

Note that $x$ is held fixed in the comparison: for each frozen $x$, the single-variable function $y \mapsto f(x, y)$ must be Lipschitz, with one constant $L$ working uniformly for all $x$ in the rectangle. The practical test carries over via the partial derivative: if $\partial f / \partial y$ is continuous on the closed rectangle $R$, then it is bounded there (a continuous function on a closed, bounded region attains a maximum), and the Mean Value Theorem in the $y$-variable gives the Lipschitz condition with
\[
  L = \max_{(x,y) \in R} \left| \frac{\partial f}{\partial y}(x, y) \right|.
\]
This is the criterion to reach for first: \emph{continuous $\partial f/\partial y$ near the initial point implies the hypotheses of the theorem hold there}.

\begin{example}
  $f(x, y) = x \sin y + y^2$ on the rectangle $|x| \leq 1$, $|y| \leq 2$. Here $\partial f/\partial y = x \cos y + 2y$, which is continuous, and on the rectangle $|x \cos y + 2y| \leq 1 + 4 = 5$. So $f$ is Lipschitz in $y$ with $L = 5$, and the theorem below will apply at any initial point inside this rectangle.
\end{example}

\section{Statement of the Theorem}

\begin{theorem}[Picard--Lindel\"of]
  Suppose $f(x, y)$ is continuous on a rectangle
  \[
    R = \{(x, y) : |x - x_0| \leq a,\ |y - y_0| \leq b\}
  \]
  and satisfies a Lipschitz condition in $y$ on $R$. Then there exists an interval $|x - x_0| \leq h$ (with $0 < h \leq a$, possibly smaller than $a$) on which the initial value problem
  \[
    y' = f(x, y), \qquad y(x_0) = y_0
  \]
  has exactly one solution.
\end{theorem}

Two features of the statement deserve emphasis.

\emph{The two conclusions have different price tags.} Continuity of $f$ alone is enough to guarantee that a solution \emph{exists} (this weaker result is Peano's existence theorem, whose proof is harder and non-constructive). Continuity alone is \emph{not} enough for uniqueness. The Lipschitz condition is what buys uniqueness---and, as a bonus, a constructive proof of existence.

\emph{The conclusion is local.} The theorem promises a solution only on some interval $|x - x_0| \leq h$, and $h$ may be strictly smaller than $a$. This is not a defect of the proof; solutions of nonlinear equations genuinely can cease to exist after finite time, as Example~\ref{ex:blowup} will show.

\section{Why Uniqueness Can Fail}
\label{sec:nonuniqueness}

Uniqueness is not a formality. Consider
\[
  y' = y^{1/3}, \qquad y(0) = 0.
\]
The constant function $y(x) \equiv 0$ is clearly a solution. But separating variables produces another: from $y^{-1/3}\,dy = dx$ we get $\tfrac{3}{2} y^{2/3} = x + C$, and the initial condition forces $C = 0$, giving
\[
  y(x) = \left( \tfrac{2}{3} x \right)^{3/2} \qquad (x \geq 0).
\]
A direct check confirms it: $y' = \tfrac{3}{2} \cdot \tfrac{2}{3} \left( \tfrac{2}{3} x \right)^{1/2} = \left( \tfrac{2}{3} x \right)^{1/2} = y^{1/3}$. Worse still, for any $c > 0$ the hybrid function
\[
  y_c(x) =
  \begin{cases}
    0, & x \leq c, \\[2pt]
    \left( \tfrac{2}{3}(x - c) \right)^{3/2}, & x > c,
  \end{cases}
\]
is also a solution: it sits at zero for a while and then ``peels off.'' Infinitely many distinct solution curves pass through the single point $(0, 0)$.

The diagnosis is exactly the failure identified in Section 2: $f(y) = y^{1/3}$ is continuous at $y = 0$ but not Lipschitz there, because its graph has a vertical tangent---the secant slopes $|y^{1/3} - 0|/|y - 0| = |y|^{-2/3}$ blow up as $y \to 0$. Peano's theorem still guarantees existence (and indeed solutions exist, too many of them); it is uniqueness that dies with the Lipschitz condition. Away from $y = 0$ the function $y^{1/3}$ is perfectly Lipschitz on any closed interval not containing the origin, so through any initial point $(x_0, y_0)$ with $y_0 \neq 0$ the solution is locally unique. The pathology lives only where the Lipschitz condition fails.

\section{The Idea of the Proof: From ODE to Integral Equation}

The proof of the theorem begins with a change of viewpoint. Integrate both sides of $y'(t) = f(t, y(t))$ from $x_0$ to $x$ and use the initial condition:
\begin{equation}
  y(x) = y_0 + \int_{x_0}^{x} f\bigl(t, y(t)\bigr)\, dt.
  \label{eq:volterra}
\end{equation}
By the Fundamental Theorem of Calculus, a continuous function $y$ satisfies the integral equation~\eqref{eq:volterra} if and only if it is differentiable and satisfies the initial value problem~\eqref{eq:ivp}. (Note how the initial condition, which in the ODE formulation is a separate side condition, has been absorbed into the equation itself as the constant term $y_0$.) Equation~\eqref{eq:volterra} is a \emph{Volterra integral equation}: the unknown function appears both outside and inside an integral whose upper limit is the variable $x$.

Why is this reformulation progress? Because it converts the problem into a \emph{fixed-point problem}. Define an operator $T$ that takes a continuous function $\varphi$ and returns a new function:
\[
  (T\varphi)(x) = y_0 + \int_{x_0}^{x} f\bigl(t, \varphi(t)\bigr)\, dt.
\]
Then $y$ solves the IVP exactly when $T y = y$, i.e., when $y$ is a fixed point of $T$. Fixed points suggest a strategy familiar from elsewhere in mathematics: \emph{iterate}. Start with a crude guess and apply $T$ repeatedly, hoping the sequence settles down.

\section{Picard Iteration}

The iteration scheme, called \emph{Picard iteration} (or the method of successive approximations), is
\[
  y_0(x) = y_0 \quad \text{(the constant function)}, \qquad
  y_{n+1}(x) = y_0 + \int_{x_0}^{x} f\bigl(t, y_n(t)\bigr)\, dt.
\]
Each step feeds the current approximation into the right-hand side and integrates to produce a better one.

The role of the Lipschitz condition is to guarantee that this process converges. A short computation using $|f(t, y_{n}) - f(t, y_{n-1})| \leq L |y_n - y_{n-1}|$ shows that on an interval of length $h$ around $x_0$,
\[
  \max_x |y_{n+1}(x) - y_n(x)| \;\leq\; Lh \cdot \max_x |y_n(x) - y_{n-1}(x)|.
\]
If $h$ is chosen small enough that $Lh < 1$, each iteration shrinks the gap between successive approximations by a fixed factor---the operator $T$ is a \emph{contraction}---and the sequence $(y_n)$ converges (by comparison with a geometric series) to a limit function that satisfies the integral equation, hence the IVP. Uniqueness follows from the same estimate: if $y$ and $\tilde{y}$ were two solutions, applying the Lipschitz bound to their difference forces $\max |y - \tilde{y}| \leq Lh \cdot \max |y - \tilde{y}|$, which with $Lh < 1$ is only possible if the difference is identically zero. (A refinement of the estimate removes the restriction $Lh<1$ and shows convergence on the whole interval where the hypotheses hold, but the contraction picture captures the essential mechanism.)

This is an instance of a far-reaching principle, the \emph{Banach fixed-point theorem}: a contraction mapping on a suitable space has exactly one fixed point, and iterating from any starting guess converges to it. Here the ``points'' of the space are continuous functions, and the ``distance'' between two of them is the maximum of their pointwise difference.

\section{A Worked Example: Watching the Iterates Converge}

Take the simplest interesting IVP, whose solution we know in advance:
\[
  y' = y, \qquad y(0) = 1, \qquad \text{exact solution } y = e^x.
\]
The integral-equation form is $y(x) = 1 + \int_0^x y(t)\, dt$, and the iteration reads $y_{n+1}(x) = 1 + \int_0^x y_n(t)\, dt$.

Starting from $y_0(x) = 1$:
\begin{align*}
  y_1(x) &= 1 + \int_0^x 1 \, dt = 1 + x, \\[4pt]
  y_2(x) &= 1 + \int_0^x (1 + t)\, dt = 1 + x + \frac{x^2}{2}, \\[4pt]
  y_3(x) &= 1 + \int_0^x \Bigl( 1 + t + \frac{t^2}{2} \Bigr) dt = 1 + x + \frac{x^2}{2} + \frac{x^3}{6}, \\[4pt]
  y_4(x) &= 1 + x + \frac{x^2}{2} + \frac{x^3}{6} + \frac{x^4}{24}.
\end{align*}
The pattern is unmistakable. Integrating $t^k/k!$ term by term produces $x^{k+1}/(k+1)!$, so each pass through the iteration appends exactly one more term of a familiar series:
\[
  y_n(x) = \sum_{k=0}^{n} \frac{x^k}{k!}
  \;\longrightarrow\;
  \sum_{k=0}^{\infty} \frac{x^k}{k!} = e^x.
\]
The abstract convergence guaranteed by the theorem is here visible in complete detail: the Picard iterates are precisely the Taylor polynomials of the solution, converging for every $x$. The global convergence reflects the fact that $f(y) = y$ is globally Lipschitz with $L = 1$; no shrinking of the interval is needed.

\section{A Nonlinear Example and the Meaning of ``Local''}

\begin{example}[Blow-up in finite time]
  \label{ex:blowup}
  Consider
  \[
    y' = y^2, \qquad y(0) = 1.
  \]
  Here $f(y) = y^2$ is locally but not globally Lipschitz (its derivative $2y$ is unbounded), so the theorem promises only a local solution. Separation of variables reveals why this caution is warranted: the exact solution is
  \[
    y(x) = \frac{1}{1 - x},
  \]
  which exists on $(-\infty, 1)$ and satisfies $y(x) \to \infty$ as $x \to 1^{-}$. The solution does not merely become hard to compute at $x = 1$; it ceases to exist. No theorem could honestly promise a solution on all of $\R$, because there isn't one.

  The first Picard iterates for this problem are
  \[
    y_0 = 1, \qquad
    y_1 = 1 + x, \qquad
    y_2 = 1 + x + x^2 + \frac{x^3}{3}, \qquad \ldots
  \]
  and one can check that their initial terms build up the geometric series $1 + x + x^2 + x^3 + \cdots = \frac{1}{1-x}$, whose radius of convergence---exactly $1$---quietly encodes the blow-up time.
\end{example}

The general lesson: for nonlinear equations, ``a solution exists'' and ``a solution exists for all $x$'' are very different claims. Picard--Lindel\"of delivers the first; the second requires additional structure (for instance, a global Lipschitz condition, which linear equations $y' = p(x)y + q(x)$ with continuous coefficients enjoy on any closed interval---this is why linear IVPs never blow up in the interior of the interval where their coefficients are continuous).

\section{Significance}

Why does this theorem occupy a foundational place in the subject?

\emph{It licenses the solving of equations.} Every solution technique implicitly assumes there is something to find and that finding one candidate finishes the job. The theorem is the warrant for both assumptions. When you solve a first-order linear equation with an integrating factor and write down ``the'' solution, uniqueness is what justifies the definite article.

\emph{It underwrites the geometric picture.} The qualitative theory of ODEs pictures the plane filled with solution curves of $y' = f(x,y)$, one through each point, never crossing. That picture \emph{is} the Picard--Lindel\"of theorem: uniqueness is exactly the statement that two distinct solution curves cannot pass through the same point. Where the Lipschitz condition fails, as in Section~\ref{sec:nonuniqueness}, curves really do merge and split, and the tidy picture breaks down.

\emph{It justifies numerical methods.} Euler's method and its refinements produce approximations to ``the solution''---language that only makes sense if a unique solution exists. Moreover, the same Lipschitz constant that drives the proof reappears in the error analysis of numerical schemes, controlling how fast nearby approximate solutions can drift apart.

\emph{Its proof method radiates outward.} The strategy---recast the problem as a fixed-point equation for an integral operator, then iterate---is a template used throughout analysis. The same contraction-mapping argument proves the inverse function theorem, solves Volterra and Fredholm integral equations via successive approximations (resolvent kernels are built by exactly this kind of iteration), and establishes local existence results for certain partial differential equations. The theorem generalizes directly to systems $\mathbf{y}' = \mathbf{f}(x, \mathbf{y})$ of first-order equations, and since any higher-order ODE can be rewritten as a first-order system, existence and uniqueness for those follows too.

\section{Self-Check Questions}

\begin{enumerate}[leftmargin=*, itemsep=8pt]

  \item State informally what the two conclusions of the Picard--Lindel\"of theorem are, and which hypothesis is responsible for each.
    \begin{answer}
      Existence of a solution near the initial point, and its uniqueness. Continuity of $f$ suffices for existence (Peano); the Lipschitz condition in $y$ is what secures uniqueness (and gives a constructive proof of existence via iteration).
    \end{answer}

  \item Give the definition of ``$g$ is Lipschitz on $I$'' and its geometric meaning.
    \begin{answer}
      There is a constant $L$ with $|g(u) - g(v)| \leq L|u-v|$ for all $u, v \in I$. Geometrically: every secant slope of the graph is bounded by $L$; the graph never becomes arbitrarily steep.
    \end{answer}

  \item Is $g(u) = |u|$ Lipschitz on $\R$? Is it differentiable? What does this tell you about the relationship between the two notions?
    \begin{answer}
      Yes, Lipschitz with $L = 1$ (triangle inequality); no, not differentiable at $0$. So Lipschitz does not imply differentiable---corners are allowed, vertical tangents are not.
    \end{answer}

  \item What is the standard practical test for verifying that $f(x,y)$ is Lipschitz in $y$ on a rectangle, and why does it work?
    \begin{answer}
      Check that $\partial f / \partial y$ is continuous on the closed rectangle; it is then bounded there by some $L$, and the Mean Value Theorem in the $y$-variable converts the bound $|\partial f/\partial y| \leq L$ into $|f(x,y_1) - f(x,y_2)| \leq L|y_1 - y_2|$.
    \end{answer}

  \item The IVP $y' = y^{1/3}$, $y(0) = 0$ has infinitely many solutions. Which hypothesis of the theorem fails, where, and why?
    \begin{answer}
      The Lipschitz condition in $y$ fails at $y = 0$: secant slopes of $y^{1/3}$ through the origin equal $|y|^{-2/3}$, which is unbounded as $y \to 0$ (vertical tangent). Continuity holds, so existence survives; uniqueness does not.
    \end{answer}

  \item Write down the Volterra integral equation equivalent to $y' = f(x,y)$, $y(x_0) = y_0$, and explain where the initial condition went.
    \begin{answer}
      $y(x) = y_0 + \int_{x_0}^{x} f(t, y(t))\, dt$. The initial condition is absorbed into the equation as the constant term $y_0$ (set $x = x_0$: the integral vanishes and $y(x_0) = y_0$ automatically).
    \end{answer}

  \item Describe the Picard iteration scheme and the role the Lipschitz condition plays in its convergence.
    \begin{answer}
      $y_0(x) \equiv y_0$ and $y_{n+1}(x) = y_0 + \int_{x_0}^x f(t, y_n(t))\, dt$. The Lipschitz bound shows each iteration multiplies the maximal gap between successive approximations by at most $Lh$; choosing $h$ with $Lh < 1$ makes the operator a contraction, so the iterates converge and the fixed point (solution) is unique.
    \end{answer}

  \item For $y' = y$, $y(0) = 1$, what are the Picard iterates $y_1, y_2, y_3$, and what do they converge to?
    \begin{answer}
      $y_1 = 1 + x$, \ $y_2 = 1 + x + x^2/2$, \ $y_3 = 1 + x + x^2/2 + x^3/6$: the Taylor polynomials of $e^x$, converging to $e^x$ for every $x$.
    \end{answer}

  \item The theorem is ``local'': it guarantees a solution only on some interval $|x - x_0| \leq h$. Give a concrete example showing that a global guarantee would be false.
    \begin{answer}
      $y' = y^2$, $y(0) = 1$ has the solution $y = 1/(1-x)$, which blows up as $x \to 1^-$. No solution exists past $x = 1$, so no theorem can promise one; $f(y) = y^2$ is only locally Lipschitz.
    \end{answer}

  \item Why do solutions of linear equations $y' = p(x)y + q(x)$ with continuous coefficients never exhibit this finite-time blow-up on the interval of continuity?
    \begin{answer}
      There $f(x,y) = p(x)y + q(x)$ satisfies $|f(x,y_1) - f(x,y_2)| = |p(x)||y_1 - y_2| \leq L|y_1-y_2|$ with $L = \max |p|$ on any closed subinterval---a Lipschitz condition holding for \emph{all} $y$, not just bounded $y$. This global-in-$y$ Lipschitz property yields existence and uniqueness on the whole interval.
    \end{answer}

  \item Uniqueness has a geometric interpretation in terms of solution curves in the plane. What is it?
    \begin{answer}
      Two distinct solution curves of $y' = f(x,y)$ can never pass through the same point (never cross or touch) in a region where the hypotheses hold. The plane is neatly partitioned into non-intersecting solution curves.
    \end{answer}

  \item (Slightly harder.) In the proof, why does the estimate $\max|y - \tilde{y}| \leq Lh \cdot \max|y - \tilde{y}|$ with $Lh < 1$ force two solutions $y$ and $\tilde{y}$ to coincide?
    \begin{answer}
      Writing $M = \max|y - \tilde{y}| \geq 0$, the inequality reads $M \leq (Lh) M$ with $Lh < 1$. If $M > 0$ we could divide by $M$ to get $1 \leq Lh$, a contradiction. Hence $M = 0$, i.e., $y = \tilde{y}$ everywhere on the interval.
    \end{answer}

\end{enumerate}

\end{document}
